\section{Normalizer \(899\), Pell unit, and arithmetic structure of the mode \(30\)}

The normalizer has the factorization arithmetic

\begin{equation}
899=30^2-1=29\cdot31.
\label{mab:rm:eq:899}
\end{equation}

The same mode \(30\) that generates sectors \(90\) and \(120\) simultaneously
determines the twin-prime pair \(29,31\). Moreover,

\begin{equation}
\varepsilon_{30}=30+\sqrt{899},\qquad
\varepsilon_{30}^{-1}=30-\sqrt{899},\qquad
\log\varepsilon_{30}=\arcosh30.
\label{mab:rm:eq:pell30}
\end{equation}

This Pell unit determines the hyperbolic scale associated with the
normalizer and relates the dodecaphase phase, mode \(30\), twin
primes, and hyperbolic geometry. This relation does not identify the
constitutive operator of the vacuum with the doubling renormalizer; such an
identification would require an explicit intertwiner.

The unit acts integrally through
\begin{equation}
M_{30}=\begin{pmatrix}30&899\\1&30\end{pmatrix},
\qquad \det M_{30}=1.
\label{mab:rm:eq:pell-matrix30}
\end{equation}
On \(\Z[\sqrt{899}]\), this matrix represents multiplication by
\(30+\sqrt{899}\). Its eigenvalues are \(\varepsilon_{30}^{\pm1}\); therefore
\(\log\varepsilon_{30}\) is the translation length of the associated hyperbolic
transformation. The same central integer \(30\) reappears in the vacuum
because \(90=3\cdot30\) and \(120=4\cdot30\), but the two operators continue to
have distinct domains. The demonstrated confluence concerns the mode and its
completions; identity of the dynamics is not postulated.

